% !PRINTED 589
\weberpage{589}
%<<CONT

Das Produkt dieser Wurzeln ist also, da $\boldsymbol{s}^m$ den Koeffizienten 1
hat, gleich dem negativen unabhängigen Glied, d.\,h. es ist:
\weq{22}{\varepsilon\,\frac{\boldsymbol{s}(\mu u)}{\boldsymbol{s}(u)} = \overset{\varrho}{\prod} \boldsymbol{s}\left(u + \frac{2\varrho}{\mu}\right),}
und daraus für $u = 0$ wie oben:
\weq{23}{\varepsilon\mu = \overset{\varrho}{\prod} \boldsymbol{s}\left(\frac{2\varrho}{\mu}\right).}

Es sei jetzt $\eta$ eine gebrochene Zahl in $\Omega$ mit dem ungeraden
Idealnenner $\mathfrak{a}$; dann folgt aus (22), daß $\boldsymbol{s}(\mu\eta)/\boldsymbol{s}(\eta)$ eine ganze
Zahl ist, wenn $\mu$ eine beliebige Zahl (6) in $\Omega$ ist. Nehmen wir
$\mu$ relativ prim zu $\mathfrak{a}$ und setzen $\mu'\mu \equiv 1$ (mod $\mathfrak{a}$), so können wir
$\eta$ durch $\mu'\eta$ ersetzen und finden, daß auch $\boldsymbol{s}(\eta)/\boldsymbol{s}(\mu'\eta)$
eine ganze Zahl ist, worin $\mu'$ ebenso beliebig ist wie $\mu$. Folglich ist
auch $\boldsymbol{s}(\eta)/\boldsymbol{s}(\mu\eta)$ eine ganze Zahl und mithin eine \gesperrt{Einheit}.

Wir haben also den Satz:

3. \gesperrt{Ist $\eta$ eine Zahl in $\Omega$ mit ungeradem Idealnenner,
und durchläuft $\mu$ eine Reihe relativer Primzahlen
zu $\mathfrak{a}$, so sind die Zahlen}
\[
\boldsymbol{s}(\mu\eta)
\]
\gesperrt{miteinander assoziiert.}

Noch eine weitere Folgerung ergibt sich daraus, wenn man
in (22) $u = 2\varrho_1/\mu_1$ setzt:
\weq{24}{\frac{\boldsymbol{s}\left(\dfrac{2\mu\varrho_1}{\mu_1}\right)}{\boldsymbol{s}\left(\dfrac{2\varrho_1}{\mu_1}\right)} = \overset{\varrho}{\prod} \boldsymbol{s}\left(\frac{2\varrho_1}{\mu_1} + \frac{2\varrho}{\mu}\right).}

Sind nun $\mu$ und $\varrho_1$ relativ prim zu $\mu_1$, so ist die linke Seite
nach 3. eine Einheit, und folglich müssen auch alle Faktoren der
rechten Seite
\weq{25}{\boldsymbol{s}\left(\frac{2\varrho_1}{\mu_1} + \frac{2\varrho}{\mu_2}\right),}
% [source sic: equation (25) prints \mu_2, whereas equation (24) and the context use \mu]
die ja ganze Zahlen sind, Einheiten sein.

Auf diese Form läßt sich aber jedes $\boldsymbol{s}(\eta)$ bringen, wenn
\[
\eta = \frac{\nu}{\mu}
\]
eine gebrochene Zahl in $\Omega$ ist, in deren Nenner zwei verschiedene
Primideale aufgehen, falls dieser Nenner ungerade ist. Denn man
kann in diesem Falle den Nenner $\mu$ von $\eta$ in zwei Ideale $\mathfrak{a}_1\mathfrak{a}_2$
%CONT>>
